\documentclass[11pt]{article}

\usepackage[margin=1in]{geometry}
\usepackage{amsmath,amssymb,amsthm}
\usepackage{enumitem}
\usepackage[colorlinks=true,linkcolor=blue,urlcolor=blue]{hyperref}

\newcommand{\MSO}{\ensuremath{\mathrm{MSO}}}
\newcommand{\MSOtwo}{\ensuremath{\mathrm{MSO}_2}}
\newcommand{\Bags}{\ensuremath{\mathrm{BAGS}}}
\newcommand{\adh}{\ensuremath{\mathrm{adh}}}
\newcommand{\cone}{\ensuremath{\mathrm{cone}}}

\title{Mathematical Issues and Gaps in\\
  \texttt{lecture\_note\_expanded.tex}}
\author{}
\date{Review of the version dated June 19, 2026}

\begin{document}
\maketitle

\section*{Scope and summary}

This report reviews the mathematical argument in
\texttt{lecture\_note\_expanded.tex}.  It does not attempt a Lean design review.
Line numbers below refer to that source file.  The most serious findings are:
\begin{enumerate}[label=\textbf{C\arabic*.},leftmargin=*]
  \item the claimed reduction of \MSOtwo{} to ordinary graph \MSO{} on the
    uncoloured incidence graph is false;
  \item the definition of a legal tree omits the requirement that the root
    boundary be empty, so the decoding and inverse lemmas are false as stated;
  \item the cone-gluing identity is not an equality of the rooted, labelled
    graphs that were defined; and
  \item the running-time statement cannot hold for the stated input format,
    since it omits both the formula length and the size of the supplied
    decomposition.
\end{enumerate}
The remaining findings are mainly missing lemmas or underspecified standard
results.  They appear repairable, but several of them are essential links in
the proof rather than cosmetic details.

\section{Critical mathematical problems}

\subsection{The incidence-graph reduction does not define the two sorts}

\paragraph{Location.} Lines 113--137, especially lines 125--135.

The note proposes, in the uncoloured incidence graph \(I(G)\),
\[
  \operatorname{edg}(z) := \exists x\,\operatorname{inc}(x,z),
  \qquad
  \operatorname{vtx}(z) := \neg\operatorname{edg}(z),
\]
and then replaces incidence by ordinary adjacency.  After that replacement,
\(\operatorname{edg}(z)\) merely says that \(z\) has a neighbour.  It holds for
every nonisolated original vertex as well as for every edge-vertex.  For
example, if \(G\) consists of one edge and its two endpoints, then every one of
the three vertices of \(I(G)\) satisfies this formula, while none satisfies the
proposed vertex formula.

More conceptually, forgetting the predicates naming the two sides of the
incidence structure loses precisely the sort information needed to interpret
typed vertex and edge variables.  Consequently the displayed translation is
not truth-preserving, and Corollary~4 is not proved.

\paragraph{Repair.} Use an incidence structure with unary predicates \(V\) and
\(E\), and prove Courcelle's theorem for a fixed finite relational vocabulary
that permits unary predicates.  This only changes the finite alphabet by
recording the unary type of each bag element.  Alternatively, attach bounded
gadgets that mark the two sorts and prove that this gadget encoding preserves a
bounded treewidth bound.  In either repair, typed first-order and monadic
quantifiers must be explicitly relativised to the appropriate unary predicate.

\subsection{A legal tree may have a nonempty root boundary}

\paragraph{Location.} Definition 6 and Lemma 7--Lemma 9, lines 248--383.

Legality checks compatibility only across parent--child edges.  It imposes no
condition on the root letter.  However, the root adhesion of every encoded
decomposition is empty.  The proof of Lemma~8 simply asserts that the root has
empty external root because it has no parent (lines 339--340), even though the
letter at the root may specify a nonempty root set.

There is an immediate counterexample: take a one-node tree whose letter is
\((H,R)\) with \(R\neq\varnothing\).  It is legal vacuously under Definition~6,
but it cannot be the encoding of any decomposition, since the root adhesion of
an encoding is empty.  Re-encoding its proposed decoding changes \(R\) to the
empty set.  Thus Lemma~9 is false as stated.

\paragraph{Repair.} Add the condition \(R=\varnothing\) at the root to the
definition of legality and to \(\varphi_{\mathrm{legal}}\).  The finite formula
can require that every root carry a letter in
\(\{(H,R):R=\varnothing\}\).

\subsection{The cone-gluing equation has the wrong external labelling}

\paragraph{Location.} Definition 5 and Lemma 6, lines 172--230.

The node graph \(H_t\) is labelled on all of \(\beta(t)\), whereas the cone graph
at \(t\) is labelled only on \(\adh(t)\).  By definition, gluing inherits the root
and the labelling from its first argument.  Hence the right-hand side of
\[
  (H_t,R_t,\iota_t)\oplus(G_1,R_1,\iota_1)
    \oplus(G_2,R_2,\iota_2)
\]
retains labels on every vertex of \(\beta(t)\), but the left-hand cone graph has
labels only on \(R_t=\adh(t)\).  The displayed equality in Lemma~6 is therefore
false as an equality of the triples defined in Definition~5.

There is also a structural equality issue: disjoint union followed by quotient
identifications normally produces a graph canonically isomorphic to the cone
graph, not literally equal to it.

\paragraph{Repair.} After both gluings, explicitly restrict the labelling to
\(R_t\), just as the decoding construction later says to do.  State the result
as a label-preserving rooted-graph isomorphism, or define a canonical quotient
model for which literal equality is meaningful.

\subsection{The stated running time ignores part of the input}

\paragraph{Location.} Theorem 1 and its proof, lines 71--78 and 806--831.

The theorem takes both a formula \(\varphi\) and a decomposition \(T\) as input,
but claims time
\[
  f(\omega,\operatorname{qr}(\varphi))\,|V(G)|.
\]
This cannot be the total running time for the stated input representation.

First, formulas of fixed quantifier rank can have arbitrarily large syntactic
length, and an algorithm must at least read the input formula.  The proof only
shows that the translated sentence and automaton depend on \(\omega\) and on the
whole formula; it does not justify replacing dependence on \(\varphi\) by
dependence only on its quantifier rank.

Second, a supplied decomposition can have arbitrarily many redundant or
duplicate nodes even when \(G\) is fixed.  Constructing its labelled encoding
takes \(\Theta(|N(T)|)\) time, as the proof itself observes.  A later appeal to a
``usual linear-size normalization'' cannot remove the cost of reading an
arbitrarily large input decomposition.

\paragraph{Repair.} A standard precise statement is either
\[
  f(\omega,|\varphi|)\bigl(|V(G)|+|N(T)|\bigr)
\]
up to the chosen input encoding, or a theorem in which \(\varphi\) is fixed and
the supplied decomposition is assumed to be in a normalized form with
\(|N(T)|=O(|V(G)|)\).  If quantifier rank is intentionally used as the logical
parameter, the compilation/preprocessing cost for reading \(\varphi\) must be
stated separately and a quantifier-rank bound for the automaton construction
must be proved.

\section{Substantial proof gaps}

\subsection{The bag-injective colouring proof is not valid as written}

\paragraph{Location.} Lines 154--168.

The greedy paragraph chooses an arbitrary bag containing the current vertex
and counts previously coloured vertices only in that bag.  It does not show
that all previously coloured vertices whose bag-subtrees intersect the current
vertex's bag-subtree occur together in the chosen bag.  Thus the asserted
greedy step is not established.  The phrase ``intersection graph of the bags''
is also inaccurate: the relevant graph has one vertex for each graph vertex,
and two such vertices conflict when their connected bag-subtrees intersect.

The existence claim is true.  A complete proof can use the facts that
intersection graphs of subtrees of a tree are chordal, subtrees of a tree have
the Helly property, and hence every clique of the conflict graph occurs in one
bag and has size at most \(\omega+1\).  Perfectness of chordal graphs then gives
the required colouring.  For the algorithmic theorem, one must additionally
give an effective ordering/colouring procedure rather than merely an existence
argument.

\subsection{Cone gluing needs intersection and edge-localisation lemmas}

\paragraph{Location.} Lines 210--230.

The proof asserts without proof that gluing identifies exactly all shared
vertices of the three pieces and that every edge of \(G[\cone(t)]\) occurs in a
bag in the subtree rooted at \(t\).  Both facts follow from the running-
intersection axiom, but neither follows merely from the definitions.

A complete proof should establish at least
\[
\begin{aligned}
  \cone(t_i)\cap\beta(t)&=\adh(t_i),\\
  \cone(t_1)\cap\cone(t_2)&\subseteq\beta(t),
\end{aligned}
\]
and then show that any edge whose endpoints lie in \(\cone(t)\) is recorded in
the parent bag or in one child cone.  The latter requires following the
connected occurrence subtrees of both endpoints to a bag covering the edge.
These lemmas are also needed later for decoding correctness.

\subsection{The decoding needs a precise quotient and a proof of connectedness}

\paragraph{Location.} Lines 294--359.

Phrases such as ``a vertex represented by \(v(t,i)\)'' presuppose an equivalence
relation generated by the gluing identifications, but that relation is not
defined.  Consequently it is not yet precise when two representatives denote
the same decoded vertex, or why a representative has a well-defined colour.

The connectedness argument in lines 352--358 also skips its central step.  It
claims that a representative can be identified with one outside a selected
component only through \(R_t\), without proving this from the recursive gluings.
A clean repair is to define decoded vertices as equivalence classes of pairs
\((t,i)\), where identifications are generated across a parent--child edge
exactly for colours in the child root.  One then proves that each equivalence
class has a connected set of tree nodes, either directly from paths in the tree
or by induction through the bottom-up gluing.  This simultaneously proves the
running-intersection axiom for the decoded bags.

\subsection{The re-encoding argument omits decisive coherence lemmas}

\paragraph{Location.} Lemma 9, especially lines 378--383.

The proof asserts that the decoded node graph is exactly the original letter.
That does not follow immediately from the definition of the decoded graph,
which is the union of edges contributed at all tree nodes.  One must rule out
the possibility that an edge contributed elsewhere becomes an additional edge
between two vertices of the decoded bag at \(t\).  Local legality should imply
the needed edge coherence, but that implication is not proved.

Likewise, one must prove that the intersection of the decoded child and parent
bags is exactly the root set recorded by the child letter.  Merely knowing that
the recorded root is contained in both bags gives only one inclusion.

The missing statements can be organized as follows:
\begin{enumerate}[label=(\alph*),leftmargin=*]
  \item two representatives \((s,i)\) and \((t,i)\) are equivalent exactly when
    the unique path between \(s\) and \(t\) carries colour \(i\) through every
    intervening child root;
  \item for a child \(x\) of \(y\),
    \(\beta(x)\cap\beta(y)\) is precisely the set represented by \(R_x\);
  \item the subgraph of the final decoded graph induced on \(\beta(t)\) is
    exactly \(H_t\).
\end{enumerate}
With these lemmas and the empty-root correction, encoding after decoding really
does return the original labelled tree.

\subsection{The formula-translation correctness proof is only announced}

\paragraph{Location.} Lines 747--802.

The note says that a structural induction gives correctness, but does not state
or prove the induction invariant.  In particular, the reverse direction of an
existential-quantifier step uses the fact that every tuple satisfying the
vertex or set guard decodes to a graph object, while the forward direction uses
surjectivity of the defining-tuple representation.  Those are the main logical
steps and should be written out.

For a syntactic translation, all auxiliary variables in the formulas for
connectedness, top nodes, and defining tuples must also be chosen fresh so that
substitution does not capture variables.  This is usually handled silently by
working modulo alpha-renaming, but the convention should be stated.  The
correctness theorem should quantify over graph assignments and their
corresponding tuple assignments, and then prove the atomic, Boolean, individual-
quantifier, and set-quantifier cases separately.

\subsection{The exact effective tree-automata theorem is not stated}

\paragraph{Location.} Lines 823--831.

The equivalence between regularity and \MSO{} definability is cited only as an
existence theorem.  To obtain an algorithm, the proof needs the effective
direction: from the particular tree \MSO{} sentence and finite alphabet, one
can compute a finite bottom-up automaton.  It should also specify whether the
trees are full binary ranked trees, ordered trees of outdegree at most two, or
trees padded by null leaves.  The automaton model and the logical vocabulary
must match this choice.

This is a legitimate external theorem, but the exact version required by the
proof should be isolated as a theorem with a citation.  If deterministic
automata are desired, include the standard determinisation step.

\section{Definitions and boundary cases requiring clarification}

\subsection{Unary nodes are omitted from the recursive construction}

The input is called a binary tree-decomposition, which often means that every
node has at most two children.  Nevertheless, the cone-gluing lemma and decoding
recursion discuss only leaves and nodes with two children.  A node with exactly
one child is not covered.  Either ``binary'' must mean full binary and a
normalisation to that form must be given, or the one-child gluing case must be
added.  This choice must agree with the tree-automaton model.

\subsection{The ambient class of tree structures is implicit}

The formulas use uniqueness of parents, a unique root, acyclicity, and exactly
one letter per node.  These are valid because the text informally evaluates
only on labelled rooted trees.  If satisfaction is instead taken over arbitrary
structures for the vocabulary
\(\{\operatorname{child}_1,\operatorname{child}_2,P_a\}\), then tree axioms and
the exactly-one-letter axiom are missing.  The cleanest mathematical statement
is to define a \(\Sigma\)-tree as a finite rooted ordered tree together with a
labelling function \(\rho:N(T)\to\Sigma\), and restrict all logical languages and
automata to that class.

\subsection{Rooted-graph notation and equality need one convention}

Definition~5 uses labels in \([k]\), while the later colour set is
\(\{0,\ldots,\omega\}\); it also introduces \(\oplus_k\) and later drops the
subscript.  These are harmless only after an explicit identification of the
label types.  More importantly, all gluing statements should consistently use
either literal equality of a chosen quotient construction or isomorphism
classes of rooted labelled graphs.

\subsection{The incidence construction also needs a size bound}

Even after repairing the two sorts, \(I(G)\) has \(|V(G)|+|E(G)|\) vertices and
the displayed decomposition adds one bag per edge.  To retain a linear bound in
\(|V(G)|\), the proof should invoke the standard fact that a simple graph of
treewidth at most \(\omega\) has \(O(\omega|V(G)|)\) edges (for example via
\(\omega\)-degeneracy).  For fixed \(\omega\) this gives the required linear
size.  Without that sentence, the runtime conclusion does not follow from the
incidence construction alone.

\section{Suggested order of repair}

The dependencies suggest the following order:
\begin{enumerate}[leftmargin=*]
  \item correct the main complexity statement and decide the exact binary-tree
    representation;
  \item replace the \MSOtwo{} incidence argument by a coloured incidence
    structure (or a proved gadget reduction);
  \item repair legality by enforcing an empty root boundary;
  \item define gluing and decoded vertices precisely, preferably using an
    explicit equivalence relation on representatives;
  \item prove the path/equivalence, adhesion, edge-coherence, and cone-union
    lemmas before restating the encoding--decoding inverse theorem;
  \item replace the colouring paragraph by a complete subtree-intersection
    argument and, for the algorithmic theorem, an effective colouring method;
  \item state and prove the assignment-level translation theorem; and
  \item state the precise effective \MSO{}-to-tree-automaton theorem being used.
\end{enumerate}

After these repairs, the defining-pair and defining-tuple part of the note
(roughly Lemmas 10--19) appears mathematically sound; its main dependency is a
fully proved encoding--decoding theorem.

\end{document}
